\chapter{คลังรหัสและองค์ประกอบสกิลเว็บสมัยใหม่ Flutter และ DevTools}

ภาคผนวกนี้รวบรวมข้อมูลจำเพาะ โค้ดต้นแบบ และสถาปัตยกรรมระดับองค์กรของสกิลการพัฒนาเว็บสมัยใหม่ แอปพลิเคชันมือถือ Flutter และเครื่องมือดีบั๊กประจำเล่มที่ 4

\section{สกิล: flutter-app-architect}

\begin{teachingplan}{ข้อมูลจำเพาะและเมทาดาทาของสกิล}
\textbf{ชื่อสกิล:} \texttt{flutter-app-architect}\\
\textbf{วัตถุประสงค์:} วางโครงสร้างแอปพลิเคชัน Flutter แบบ Clean Architecture (MVVM/Provider) ปราศจากข้อผิดพลาด Static Linter\\
\textbf{มาตรฐานการประเมิน:} Multi-Screen Responsive, 0 Errors, 0 Warnings, Custom App Icon Pipeline
\end{teachingplan}

\subsection{โค้ดแกนกลาง: แม่แบบ ViewModel และ State Management}

\begin{lstlisting}[language=JavaScript, caption={Flutter ViewModel Architecture Blueprint}]
class PhysicsExperimentViewModel extends ChangeNotifier {
  double _stringLength = 1.0; // meters
  double _gravity = 9.80665;  // m/s^2
  bool _isRunning = false;

  double get stringLength => _stringLength;
  double get period => 2 * 3.14159265 * Math.sqrt(_stringLength / _gravity);
  bool get isRunning => _isRunning;

  void updateLength(double newLength) {
    if (newLength > 0.1 && newLength <= 3.0) {
      _stringLength = newLength;
      notifyListeners(); // Triggers atomic UI re-render
    }
  }

  void toggleSimulation() {
    _isRunning = !_isRunning;
    notifyListeners();
  }
}
\end{lstlisting}

\section{สกิล: modern-web-guidance}

\begin{teachingplan}{ข้อมูลจำเพาะและเมทาดาทาของสกิล}
\textbf{ชื่อสกิล:} \texttt{modern-web-guidance}\\
\textbf{วัตถุประสงค์:} กำหนดแนวปฏิบัติและมาตรฐาน API เว็บฟรอนต์เอนด์ยุคใหม่เพื่อประสิทธิภาพสูงสุด\\
\textbf{ขอบเขตเทคโนโลยี:} CSS Container Queries, View Transitions, Zero-GC RequestAnimationFrame
\end{teachingplan}

\subsection{โค้ดแกนกลาง: สถาปัตยกรรม CSS Container Queries}

\begin{lstlisting}[language=JavaScript, caption={Responsive Card Component Style}]
/* Container Definition */
.lab-card-container {
  container-type: inline-size;
  container-name: labCard;
}

/* Dynamic Layout Adaptation */
@container labCard (min-width: 450px) {
  .lab-card {
    display: grid;
    grid-template-columns: 1fr 1fr;
    gap: 1.5rem;
    padding: 1.25rem;
  }
}
\end{lstlisting}

\section{สกิล: memory-leak-debugging}

\begin{teachingplan}{ข้อมูลจำเพาะและเมทาดาทาของสกิล}
\textbf{ชื่อสกิล:} \texttt{memory-leak-debugging}\\
\textbf{วัตถุประสงค์:} วิเคราะห์ Heap Snapshot ตรวจจับหน่วยความจำรั่วไหลและระบุ Retainer Trees\\
\textbf{ขอบเขตเทคโนโลยี:} Chrome DevTools MCP, Heap Profiler, WeakMap Management
\end{teachingplan}

\subsection{โค้ดแกนกลาง: การจัดเก็บตัวแปรอ้างอิงแบบอัตโนมัติด้วย WeakMap}

\begin{lstlisting}[language=JavaScript, caption={Leak-Proof DOM Reference Cache}]
// Prevents Detached DOM Node Leaks
const nodeMetadataCache = new WeakMap();

function registerElementData(element, metadata) {
  // If element is removed from DOM, GC reclaims metadata automatically
  nodeMetadataCache.set(element, metadata);
}
\end{lstlisting}
